\documentclass[11pt,a4paper]{article}
\usepackage[utf8]{inputenc}
\usepackage[T1]{fontenc}
\usepackage{lmodern}
\usepackage{hyperref}
\usepackage{geometry}
\usepackage{amsmath}
\usepackage{array}
\usepackage{enumitem}

\geometry{margin=2.5cm}
\setlength{\parskip}{6pt}
\setlength{\parindent}{0pt}

\hypersetup{
    colorlinks=true,
    linkcolor=blue,
    urlcolor=blue,
    citecolor=blue
}

\title{\textit{foofit}: X-ray Reflectivity Fitting Software}
\author{
  H.-G. Steinr\"uck\textsuperscript{1,*} \\[4pt]
  \small \textsuperscript{1}IHE-1, Forschungszentrum J\"ulich \& RWTH Aachen \\
  \small *\,\href{mailto:h.steinrueck@fz-juelich.de}{h.steinrueck@fz-juelich.de}
}
\date{February 2026}

\begin{document}

\maketitle

\begin{abstract}
This document serves as the manual for \textit{foofit}, a pip-installable Python library
for X-ray reflectivity (XRR) fitting. \textit{foofit} implements the Parratt recursion
formalism~\cite{Parratt1954} and a refraction-corrected master formula, using
\texttt{lmfit}~\cite{Newville2014} for optimization and Monte Carlo confidence analysis.
\end{abstract}

% ----------------------------------------------------------------
\section{Introduction}
% ----------------------------------------------------------------

\textit{foofit} is a Python library for fitting X-ray reflectivity (XRR) data. It implements
two formalisms for computing theoretical XRR curves: the recursive Parratt formalism and a
refraction-corrected master formula. Optimization is handled via the \texttt{Minimizer}
function from \texttt{lmfit}~\cite{Newville2014}, and Monte Carlo resampling~\cite{Heinrich2009}
is used for parameter confidence analysis.

The code was originally developed by Michael Klimczak and Hans-Georg Steinr\"uck at the
Crystallography and Structural Physics institute of the Friedrich-Alexander Universit\"at
Erlangen-N\"urnberg (based on \textit{pyRefFit}), and later refactored into a
pip-installable package.

Source code and issue tracker: \url{https://github.com/hgstei/foofit}

% ----------------------------------------------------------------
\section{Installation}
% ----------------------------------------------------------------

Install directly from GitHub:
\begin{verbatim}
pip install git+https://github.com/hgstei/foofit.git
\end{verbatim}

All dependencies are installed automatically:
\texttt{numpy}, \texttt{scipy}, \texttt{matplotlib}, \texttt{lmfit}, \texttt{prettytable},
\texttt{tqdm}, \texttt{corner}, \texttt{joblib}, \texttt{numdifftools}, \texttt{astropy}.

For interactive, zoomable plots in Jupyter, additionally install \texttt{ipympl}:
\begin{verbatim}
pip install ipympl
\end{verbatim}

% ----------------------------------------------------------------
\section{Quick Start}
% ----------------------------------------------------------------

\begin{verbatim}
from foofit import *

# bundled example dataset: polystyrene on silicon
si_ps_xrr = example_data

# set up lmfit Parameters
params = Parameters()
params.add('numbLayers', value=1,    vary=False)
params.add('wavelength', value=1.0,  vary=False)
params.add('I0',         value=1,    vary=False)
params.add('bkg',        value=0,    vary=False)
params.add('pre_rho',    value=0,    vary=False)
params.add('pre_beta',   value=0,    vary=False)
params.add('layer0_dd',  value=220, min=150, max=300, vary=True)
params.add('layer0_rho', value=0.25, min=0.1, max=0.4, vary=True)
params.add('layer0_sig', value=5,   min=0.5, max=15,  vary=True)
params.add('layer0_beta',value=0,    vary=False)
params.add('sub_rho',    value=0.71, vary=False)
params.add('sub_sig',    value=4,   min=0.5, max=10,  vary=True)
params.add('sub_beta',   value=0,    vary=False)

# single fit
performFit(si_ps_xrr, params, fitFunc=xrr_parratt_fit,
           method='powell', qmin=0.0, qmax=0.4, plot=2,
           outputName='my_sample', weight=2)

# Monte Carlo error analysis
params2 = loadParams('my_sample_<timestamp>_fit.fitParams',
                     lowLim=0.3, highLim=0.3)
performFit_mc(si_ps_xrr, params2, fitFunc=xrr_parratt_fit,
              method='powell', qmin=0.0, qmax=0.4,
              outputName='my_sample', weight=2, NN=250)

# corner plot of MC parameter distributions
analyze_mc('my_sample_<timestamp>_fit_mc.fitParams')
\end{verbatim}

A full worked example is provided in \texttt{example.ipynb}.

\subsection*{Interactive Plots in Jupyter}

Uncomment the appropriate line at the top of the notebook:
\begin{verbatim}
# %matplotlib widget    # JupyterLab (requires: pip install ipympl)
# %matplotlib notebook  # classic Jupyter Notebook
\end{verbatim}
\textbf{Note:} Do not add comments on the same line as a \texttt{\%matplotlib}
magic command --- IPython does not treat \texttt{\#} as a comment in magic arguments.
Put the comment on the next line instead.

% ----------------------------------------------------------------
\section{Parameters Convention}
% ----------------------------------------------------------------

All fitting functions take an \texttt{lmfit} \texttt{Parameters} object. The parameter
names listed below must be followed exactly. Layer index $n$ starts at~0 for the topmost
layer (closest to the ambient medium). Mathematical constraints between parameters can be
defined as described in the lmfit documentation~\cite{lmfit_constraints}.

\begin{center}
\renewcommand{\arraystretch}{1.3}
\begin{tabular}{lp{9.5cm}}
\hline\hline
\textbf{Parameter} & \textbf{Description} \\
\hline
\texttt{numbLayers}        & Number of layers between ambient and substrate (integer, fixed) \\
\texttt{wavelength}        & X-ray wavelength in \AA{} (fixed) \\
\texttt{I0}                & Intensity scale factor \\
\texttt{bkg}               & Background offset \\
\texttt{pre\_rho}          & Ambient medium electron density (0 for air/vacuum) \\
\texttt{pre\_beta}         & Ambient medium absorption constant (0 for air/vacuum) \\
\texttt{layer\{n\}\_dd}   & Thickness of layer $n$ in \AA{} \\
\texttt{layer\{n\}\_rho}  & Electron density of layer $n$ in e$^-$/\AA$^3$ \\
\texttt{layer\{n\}\_sig}  & Interfacial roughness of layer $n$ in \AA{} \\
\texttt{layer\{n\}\_beta} & Absorption constant of layer $n$ \\
\texttt{sub\_rho}          & Substrate electron density in e$^-$/\AA$^3$ \\
\texttt{sub\_sig}          & Substrate roughness in \AA{} \\
\texttt{sub\_beta}         & Substrate absorption constant \\
\hline\hline
\end{tabular}
\end{center}

% ----------------------------------------------------------------
\section{Functions}
% ----------------------------------------------------------------

\subsection*{\texttt{xrr\_parratt\_calc(params, qq, doConv=0)}}

Calculates XRR using the recursive Parratt formalism~\cite{Parratt1954}. Takes an
\texttt{lmfit} \texttt{Parameters} object, a q-array in \AA$^{-1}$, and optionally a
Gaussian convolution width \texttt{doConv} (sigma in \AA$^{-1}$, \texttt{0} to disable).
Returns the reflectivity array. Cross-checked against GenX~\cite{Bjorck2007}.

\subsection*{\texttt{xrr\_master\_refractionCorrected\_calc(params, qq, doConv=0)}}

Same interface as \texttt{xrr\_parratt\_calc}, but uses the refraction-corrected master formula.

\subsection*{\texttt{xrr\_parratt\_fit(params, qq, data, weight, ee, doConv=0)}}

Residual function for use with \texttt{lmfit}. Returns the weighted difference between
model (Parratt) and data. Passed as \texttt{fitFunc} to \texttt{performFit} and
\texttt{performFit\_mc}.

\subsection*{\texttt{xrr\_master\_refractionCorrected\_fit(params, qq, data, weight, ee, doConv=0)}}

Same as \texttt{xrr\_parratt\_fit} but using the refraction-corrected master formula.

\subsection*{\texttt{xrr\_eDens(params, zz)}}

Calculates the electron density profile (in e$^-$/\AA$^3$) including interfacial roughness
as a function of depth $z$ in \AA.

\subsection*{\texttt{xrr\_eDens\_zeroRoughness(params, zz)}}

Same as \texttt{xrr\_eDens} but with sharp (roughness-free) interfaces.

\subsection*{\texttt{xrr\_beta(params, zz)}}

Calculates the absorption ($\beta$) profile as a function of depth $z$ in \AA.

\subsection*{\texttt{xrr\_beta\_zeroRoughness(params, zz)}}

Same as \texttt{xrr\_beta} but with sharp (roughness-free) interfaces.

\subsection*{\texttt{performFit(\ldots)}}

\begin{verbatim}
performFit(dataFile, params, fitFunc=xrr_parratt_fit,
           method='differential_evolution',
           qmin=0, qmax=1, plot=1, report=True, save=True,
           outputName="foo", weight=0, rrfPlot=False, doConv=0)
\end{verbatim}

Performs a single fit and optionally plots and saves results.

\begin{itemize}[leftmargin=1.5em]
  \item \texttt{dataFile}: Path to data file. Column~1 must be $q$ in \AA$^{-1}$,
        column~2 must be intensity $I$. Additional columns are ignored.
  \item \texttt{fitFunc}: \texttt{xrr\_parratt\_fit} (default) or
        \texttt{xrr\_master\_refractionCorrected\_fit}.
  \item \texttt{method}: Optimization algorithm from \texttt{lmfit}.
        \texttt{'powell'} and \texttt{'differential\_evolution'} work best in practice.
  \item \texttt{qmin}, \texttt{qmax}: Fitted q-range in \AA$^{-1}$.
  \item \texttt{plot}: \texttt{1}~plots XRR and fit only; \texttt{2}~additionally saves
        all output files (see below).
  \item \texttt{report}: Print fit report to console.
  \item \texttt{save}: Save report plot to file.
  \item \texttt{outputName}: Base name for all output files.
  \item \texttt{weight}: Figure-of-merit definition.
        \texttt{0}: $(\mathrm{data}-\mathrm{model})/\mathrm{data}$;
        \texttt{1}: $\mathrm{data}-\mathrm{model}$;
        \texttt{2}: $\log(\mathrm{data})-\log(\mathrm{model})$;
        \texttt{4}: $(\mathrm{data}-\mathrm{model})/\mathrm{error}$.
  \item \texttt{rrfPlot}: If \texttt{True}, display Fresnel-normalized reflectivity
        ($R \cdot q^4$).
  \item \texttt{doConv}: Gaussian convolution width (sigma in \AA$^{-1}$) for instrumental
        resolution smearing. \texttt{0} disables convolution.
\end{itemize}

Output files saved when \texttt{plot=2} (all prefixed \texttt{outputName\_timestamp\_fit}):
\texttt{.fitParams},
\texttt{.ed} (electron density),
\texttt{.ned} (electron density, no roughness),
\texttt{.beta},
\texttt{.nbeta} (absorption, no roughness),
\texttt{.r} (reflectivity model curve).

\subsection*{\texttt{performFit\_mc(\ldots)}}

\begin{verbatim}
performFit_mc(dataFile, params, fitFunc=xrr_parratt_fit,
              method='differential_evolution',
              qmin=0, qmax=1, plot=1, report=True,
              outputName="foo", weight=0, rrfPlot=False,
              doConv=0, NN=10)
\end{verbatim}

Performs Monte Carlo error analysis~\cite{Heinrich2009} by generating \texttt{NN}
noise-perturbed datasets and fitting each independently. Fits are parallelized using
\texttt{joblib}~\cite{joblib}. Always saves output files and creates a report figure.
Similar in approach to MOTOFIT~\cite{Nelson2006}.

Parameters are identical to \texttt{performFit}, with the addition of:
\begin{itemize}[leftmargin=1.5em]
  \item \texttt{NN}: Number of Monte Carlo iterations (default~10;
        100--500 recommended for reliable statistics).
\end{itemize}

\subsection*{\texttt{loadParams(parameterFilename, lowLim=0.1, highLim=0.1)}}

Loads a \texttt{.fitParams} file saved by \texttt{performFit} back into an \texttt{lmfit}
\texttt{Parameters} object. \texttt{lowLim} and \texttt{highLim} define lower and upper
parameter bounds as a fraction of the fitted value (e.g., \texttt{0.3} gives $\pm30\,\%$
bounds), used as limits in the subsequent Monte Carlo fit.

\subsection*{\texttt{analyze\_mc(file, bins=20)}}

Generates a corner plot~\cite{ForemanMackey2016} from the \texttt{.fitParams} file saved
by \texttt{performFit\_mc}, showing marginal parameter distributions and pairwise correlations.

\subsection*{\texttt{example\_data}}

Module-level string variable containing the absolute path to the bundled example dataset
(\texttt{si\_ps.xrr}): a polystyrene thin film on a silicon substrate.

% ----------------------------------------------------------------
\section{Physical Constants}
% ----------------------------------------------------------------

The classical electron radius $r_e$ and the critical-$q$ prefactor are computed from
\texttt{astropy} constants at import time and re-exported by \texttt{from foofit import *}:

\begin{verbatim}
# r_e not directly available in astropy < 4;
# derived from Thomson cross-section: sigma_T = (8*pi/3) * r_e^2
r_e_AA = sqrt(3 * sigma_T / (8*pi)) * 1e10  # = 2.81794e-5  Angstrom
qc_factor = 4 * sqrt(r_e_AA * pi)           # = 0.037636
\end{verbatim}

% ----------------------------------------------------------------
\begin{thebibliography}{9}
% ----------------------------------------------------------------

\bibitem{Parratt1954}
  L.G. Parratt,
  Physical Review \textbf{95}, 359 (1954).

\bibitem{Newville2014}
  M. Newville et al.,
  \textit{LMFIT: Non-Linear Least-Square Minimization and Curve-Fitting for Python} (2014),
  \url{https://lmfit.github.io/lmfit-py/}.

\bibitem{lmfit_constraints}
  \textit{lmfit constraints documentation},
  \url{https://lmfit.github.io/lmfit-py/constraints.html}.

\bibitem{Bjorck2007}
  M. Bj\"orck and G. Andersson,
  Journal of Applied Crystallography \textbf{40}, 1174 (2007).

\bibitem{Heinrich2009}
  F. Heinrich, T. Ng, D.J. Vanderah, P. Shekhar, M. Mihailescu, H. Nanda, and M. L\"osche,
  Langmuir \textbf{25}, 4219 (2009).

\bibitem{Nelson2006}
  A. Nelson,
  Journal of Applied Crystallography \textbf{39}, 273 (2006).

\bibitem{joblib}
  \textit{joblib: running Python functions as pipeline jobs},
  \url{https://joblib.readthedocs.io/}.

\bibitem{ForemanMackey2016}
  D. Foreman-Mackey,
  Journal of Open Source Software \textbf{1}, 24 (2016).

\end{thebibliography}

\end{document}
